\documentclass[11pt,a4paper]{article}
\usepackage[margin=2.5cm]{geometry}
\usepackage{booktabs}
\usepackage{longtable}
\usepackage{array}
\usepackage{xcolor}
\usepackage{hyperref}
\usepackage{amsmath}
\usepackage{enumitem}
\usepackage{caption}

\definecolor{secblue}{RGB}{0,70,140}
\hypersetup{colorlinks=true,linkcolor=secblue,urlcolor=secblue}

\newcolumntype{L}[1]{>{\raggedright\arraybackslash}p{#1}}

\title{\textbf{ML-Wheat Multi-Scale Feature Documentation}\\[6pt]
\large Variables in \texttt{wheat\_multiscale\_features\_\{year\}.parquet}}
\author{Kelly Trinh}
\date{}

\begin{document}
\maketitle
\tableofcontents
\newpage

% ======================================================================
\section{Overview}

The file \texttt{wheat\_multiscale\_features\_\{year\}.parquet} contains one row per wheat-growing grid cell ($\sim$26,736 cells per year) with features extracted at three temporal resolutions:

\begin{enumerate}[leftmargin=2cm]
    \item \textbf{Weekly} --- 60 slots $\times$ 16 variables = 960 features
    \item \textbf{Monthly} --- 15 slots $\times$ 16 variables = 240 features
    \item \textbf{Crop-seasonal} --- 6 windows $\times$ 16 variables + $\sim$25 extras $\approx$ 121 features
    \item \textbf{Static} --- 9 identifiers and soil properties
\end{enumerate}

\noindent\textbf{Total: $\sim$1,320s features per grid cell per year.}

All features are derived from four SILO climate inputs (tmin, tmax, rain, radn), gridded soil moisture (s0), and static APSIM soil properties. The APSIM-Wheat documentation (v7.5 R3008) is referenced throughout with equation numbers.

\subsection{Data Sources}

\begin{tabular}{lll}
\toprule
\textbf{Source} & \textbf{Variables} & \textbf{Resolution} \\
\midrule
SILO climate NetCDF  & tmin, tmax, rain, radn & Daily, 0.05\textdegree \\
Soil moisture (s0)   & Plant-available water (mm) & Daily, gridded \\
APSIM soil XML       & PAWC, pH, mineral N, depth & Static per soil type \\
Gridded soil map     & Dominant soil ID per cell & Static per cell \\
DEWS yield NetCDF    & Wheat yield (kg/ha $\rightarrow$ t/ha) & Annual \\
\bottomrule
\end{tabular}

\subsection{Sowing Date Determination}

Sowing dates are computed per cell per year using DEWS rules (not fixed):
\begin{itemize}
    \item \textbf{Northern} (lat $> -32.24$, lon $> 140.99$): rain $\geq$ 15\,mm/3\,days AND PAW $\geq$ 30\,mm, Apr 26--Jul 31
    \item \textbf{Southern} (lat $< -32.25$, lon $> 129$): rain $\geq$ 15\,mm/3\,days, Apr 26--Jul 31
    \item \textbf{Western} (lon $\leq 129$): rain $\geq$ 15\,mm/3\,days, Apr 1--Jul 31
    \item Wet-soil delay: PAW $<$ 95\% of PAWC. Default sow Jul 31 if no criteria met.
\end{itemize}

% ======================================================================
\section{Static Features (9 columns)}

These are time-invariant identifiers and soil properties.

\begin{longtable}{L{3.8cm}L{1.8cm}L{7.5cm}}
\toprule
\textbf{Column} & \textbf{Unit} & \textbf{Description} \\
\midrule
\endhead
\texttt{lat} & \textdegree S & Latitude of grid cell \\
\texttt{lon} & \textdegree E & Longitude of grid cell \\
\texttt{sowing\_doy} & day & Day-of-year of sowing (NaN if not sown) \\
\texttt{no\_sow} & 0/1 & 1 if no sowing occurred \\
\texttt{pawc\_0\_30\_mm} & mm & Plant Available Water Capacity, 0--30\,cm depth (APSIM Eq\,70, 89--90) \\
\texttt{pawc\_0\_60\_mm} & mm & PAWC, 0--60\,cm depth \\
\texttt{ph\_0\_30} & --- & Depth-weighted mean soil pH, 0--30\,cm \\
\texttt{minN\_0\_30} & mg/kg & Mineral nitrogen (NO$_3$ + NH$_4$), 0--30\,cm (relates to Eq\,98--99) \\
\texttt{profile\_depth\_cm} & cm & Total soil profile depth (constrains root growth, Section\,9.1) \\
\midrule
\texttt{wheat\_yield} & t/ha & APSIM-simulated yield (target variable) \\
\texttt{year} & --- & Calendar year \\
\bottomrule
\end{longtable}

% ======================================================================
\section{Temporal Slot Variables (16 per slot)}
\label{sec:slotvars}

The same 16 variables are computed for every temporal slot (weekly, monthly, and crop-seasonal). Each variable name is formed as \texttt{\{variable\}\_\{slot\_name\}}, e.g.\ \texttt{rain\_sum\_w3} or \texttt{fasw\_mean\_mpre2}.

\begin{longtable}{L{3.2cm}L{1.8cm}L{8.2cm}}
\toprule
\textbf{Variable} & \textbf{Unit} & \textbf{Description and APSIM Reference} \\
\midrule
\endhead
\texttt{rain\_sum} & mm & Total rainfall. Water input for the soil water balance (Section\,11). \\
\texttt{rain\_days} & days & Number of days with rainfall $\geq$ 1\,mm. Rainfall frequency pattern. \\
\texttt{rad\_mean} & MJ/m$^2$/d & Mean daily solar radiation. Drives biomass accumulation via RUE (Eq\,13--16). \\
\texttt{rad\_sum} & MJ/m$^2$ & Cumulative radiation. Total energy available for photosynthesis. \\
\texttt{tmean} & \textdegree C & Mean daily temperature $(\text{Tmax}+\text{Tmin})/2$. General thermal environment. \\
\texttt{tmax\_max} & \textdegree C & Maximum of daily Tmax. Detects heat extremes ($>$34\textdegree C triggers senescence, Eq\,81). \\
\texttt{tmin\_min} & \textdegree C & Minimum of daily Tmin. Detects frost events ($<$0\textdegree C, Eq\,80). \\
\texttt{diurnal} & \textdegree C & Mean diurnal temperature range (Tmax $-$ Tmin). Influences VPD and vernalisation (Eq\,9). \\
\texttt{tt\_sum} & \textdegree Cd & Thermal time accumulated (Eq\,4). Cardinal temperatures: 0/26/34\textdegree C. \\
\texttt{fT\_photo} & 0--1 & Mean temperature stress factor for photosynthesis (Eq\,20, Fig\,10). Breakpoints: 0/10/25/35\textdegree C. \\
\texttt{heat\_days} & days & Count of days with Tmax $>$ 34\textdegree C. Heat senescence events (Eq\,81, Fig\,29). \\
\texttt{frost\_days} & days & Count of days with Tmin $<$ 0\textdegree C. Frost damage potential (Eq\,80). \\
\texttt{vpd\_mean} & kPa & Mean vapour pressure deficit (Eq\,88). Drives transpiration efficiency (Eq\,87). \\
\texttt{fasw\_mean} & 0--2 & Mean fraction of available soil water: $\text{SM}/\text{PAWC}$ (proxy for Eq\,70). \\
\texttt{fw\_photo} & 0--1 & Mean photosynthesis water stress: $\min(\text{FASW}/0.5, 1)$ (proxy for Eq\,95--96). \\
\texttt{n\_days} & days & Number of valid days in the slot. Handles partial weeks/months at boundaries. \\
\bottomrule
\end{longtable}

% ======================================================================
\section{Weekly Features (704 features)}

\subsection{Slot Naming Convention}

\begin{tabular}{ll}
\toprule
\textbf{Slot name} & \textbf{Meaning} \\
\midrule
\texttt{wpre4} \ldots\ \texttt{wpre1} & 4 weeks before sowing (pre-season) \\
\texttt{w0} & Week of sowing (days 0--6 from sowing) \\
\texttt{w1} \ldots\ \texttt{w39} & Weeks 1--39 after sowing \\
\bottomrule
\end{tabular}

\medskip
Each slot spans exactly 7 days relative to the per-cell sowing date. Pre-sowing weeks capture antecedent soil moisture recharge and rainfall patterns.

\textbf{Total:} 44 slots $\times$ 16 variables = 704 features.

\textbf{Example column names:}
\texttt{rain\_sum\_wpre2}, \texttt{fasw\_mean\_w5}, \texttt{heat\_days\_w20}, \texttt{tt\_sum\_w0}.

\subsection{Approximate Phenological Alignment}

\begin{tabular}{lll}
\toprule
\textbf{Weeks from sowing} & \textbf{Approx.\ stage} & \textbf{Key processes} \\
\midrule
w0--w2 & Establishment & Germination, emergence \\
w3--w15 & Vegetative & Tillering, leaf expansion, vernalisation \\
w16--w18 & Pre-anthesis & Floral initiation, stem elongation \\
w19--w28 & Grain fill & Grain number set, grain filling \\
w29--w39 & Maturation & Senescence, harvest \\
\bottomrule
\end{tabular}

\medskip
\noindent Note: the exact alignment varies by location and year because phenological progression depends on accumulated thermal time, not calendar weeks.

% ======================================================================
\section{Monthly Features (192 features)}

\subsection{Slot Naming Convention}

\begin{tabular}{ll}
\toprule
\textbf{Slot name} & \textbf{Meaning} \\
\midrule
\texttt{mpre3} \ldots\ \texttt{mpre1} & 3 calendar months before the sowing month \\
\texttt{m0} & Sowing month (only post-sowing days included) \\
\texttt{m1} \ldots\ \texttt{m8} & Months 1--8 after sowing month \\
\bottomrule
\end{tabular}

\medskip
Monthly slots are defined by calendar months relative to each cell's sowing month. For a May sowing, \texttt{mpre1} = April, \texttt{m0} = May (post-sowing days only), \texttt{m3} = August, etc.

\textbf{Total:} 12 slots $\times$ 16 variables = 192 features.

\textbf{Pre-sowing months} capture soil moisture recharge (e.g.\ autumn break rainfall in southern Australia), stored soil water, and temperature regime before the crop is established.

\textbf{Example column names:}
\texttt{rain\_sum\_mpre1}, \texttt{fw\_photo\_m4}, \texttt{rad\_sum\_m7}.

% ======================================================================
\section{Crop-Seasonal Features ($\sim$121 features)}

\subsection{APSIM Phenological Windows}

Features are aggregated over five thermal-time-based windows (same 16 variables as Section~\ref{sec:slotvars}):

\begin{tabular}{llll}
\toprule
\textbf{Window} & \textbf{Cum.\ TT range (\textdegree Cd)} & \textbf{Phenological period} & \textbf{Slot suffix} \\
\midrule
W1\_estab & 0 -- 100 & Sowing $\rightarrow$ Emergence & \texttt{\_W1\_estab} \\
W2\_veg & 100 -- 1055 & Emergence $\rightarrow$ Floral Initiation & \texttt{\_W2\_veg} \\
W3\_preAnth & 1055 -- 1175 & Floral Init.\ $\rightarrow$ Flowering & \texttt{\_W3\_preAnth} \\
W4\_grainFill & 1175 -- 1860 & Flowering $\rightarrow$ End Grain Fill & \texttt{\_W4\_grainFill} \\
W5\_matur & 1860 -- 2060 & End Grain Fill $\rightarrow$ Maturity & \texttt{\_W5\_matur} \\
\midrule
crop & 0 -- 2060 & Whole crop cycle & \texttt{\_crop} \\
\bottomrule
\end{tabular}

\medskip
\textbf{Subtotal:} 6 windows $\times$ 16 variables = 96 features.

Thermal time targets are for the Hartog cultivar (APSIM default). Cumulative thermal time is calculated using Eq\,4 with cardinal temperatures 0/26/34\textdegree C.

\subsection{Additional Crop-Seasonal Features}

\begin{longtable}{L{4.5cm}L{2.5cm}L{6.2cm}}
\toprule
\textbf{Feature} & \textbf{Windows} & \textbf{Description (APSIM reference)} \\
\midrule
\endhead
\multicolumn{3}{l}{\textit{Water stress --- expansion}} \\
\texttt{fw\_expan\_mean\_\{W\}} & W2, W3 & Mean leaf expansion stress factor. More sensitive than \texttt{fw\_photo}: threshold FASW/0.8 (Eq\,97, Fig\,34). \\
\texttt{fw\_expan\_min\_\{W\}} & W2, W3 & Minimum (worst) expansion stress in the window. \\
\midrule
\multicolumn{3}{l}{\textit{Water stress --- photosynthesis minimum}} \\
\texttt{fw\_photo\_min\_\{W\}} & W2, W3, W4 & Worst single-day photosynthesis water stress. Captures peak drought severity. \\
\midrule
\multicolumn{3}{l}{\textit{Drought event counts}} \\
\texttt{drought\_severe\_\{W\}} & W2, W3, W4 & Days with \texttt{fw\_photo} $<$ 0.5. Severe water stress events. \\
\texttt{drought\_moderate\_\{W\}} & W2, W3, W4 & Days with $0.5 \leq$ \texttt{fw\_photo} $< 0.8$. Moderate stress events. \\
\midrule
\multicolumn{3}{l}{\textit{Cumulative water stress}} \\
\texttt{cum\_water\_stress\_crop} & Whole crop & $\sum(1 - \texttt{fw\_photo})$ over all crop days. Total season water deficit. \\
\texttt{cum\_water\_stress\_\{W\}} & W3, W4 & Cumulative stress in pre-anthesis and grain fill. \\
\midrule
\multicolumn{3}{l}{\textit{Soil water dynamics}} \\
\texttt{FASW\_at\_sowing} & Sowing $\pm$5d & Mean FASW at sowing. Initial soil water status. \\
\texttt{FASW\_W4\_drying\_rate} & W4 & FASW$_{\text{1st half}}$ $-$ FASW$_{\text{2nd half}}$ of grain fill. Soil water depletion rate during grain filling. \\
\midrule
\multicolumn{3}{l}{\textit{Developmental factors}} \\
\texttt{vern\_sum\_W2\_veg} & W2 & Cumulative vernalisation (Eq\,9--12). Drives phenological progress from emergence to floral initiation. \\
\texttt{fD\_mean\_W2\_veg} & W2 & Mean photoperiod factor (Eq\,8). Short days slow development. \\
\midrule
\multicolumn{3}{l}{\textit{Grain fill factors}} \\
\texttt{hg\_mean\_W4} & W4 & Mean grain fill temperature factor (Eq\,49, Fig\,17). Optimal at 18--25\textdegree C. \\
\texttt{heat\_sen\_W4} & W4 & Mean heat senescence factor. Non-zero when Tmax $>$ 34\textdegree C (Eq\,81). \\
\midrule
\multicolumn{3}{l}{\textit{Transpiration and total thermal time}} \\
\texttt{VPD\_mean\_crop} & Whole crop & Mean VPD (Eq\,88). Controls transpiration efficiency. \\
\texttt{TE\_mean\_crop} & Whole crop & Mean transpiration efficiency proxy: $0.006/\text{VPD}$ (Eq\,87). \\
\texttt{TT\_cum\_crop} & Whole crop & Total thermal time accumulated. Indicates season length. \\
\bottomrule
\end{longtable}

% ======================================================================
\section{APSIM Equation Summary}

Key equations from the APSIM-Wheat documentation used in feature computation:

\begin{align}
\text{Crown temperature:} \quad T_c &= \frac{T_{c,\max} + T_{c,\min}}{2} \tag{Eq\,3} \\[4pt]
\text{Thermal time:} \quad \Delta TT &= \begin{cases} T_c & 0 < T_c \leq 26 \\ \frac{26}{8}(34-T_c) & 26 < T_c \leq 34 \\ 0 & \text{otherwise} \end{cases} \tag{Eq\,4} \\[4pt]
\text{Photoperiod factor:} \quad f_D &= 1 - 0.002 R_p (20 - L_p)^2 \tag{Eq\,8} \\[4pt]
\text{Vernalisation:} \quad \Delta V &= \min\!\Big(1.4 - 0.0778 T_c,\; 0.5 + \frac{13.44 T_c}{(T_{\max}-T_{\min}+3)^2}\Big) \tag{Eq\,9} \\[4pt]
\text{Photosynthesis stress:} \quad f_{T,\text{photo}} &= h_{T,\text{photo}}\!\Big(\frac{T_{\max}+T_{\min}}{2}\Big) \tag{Eq\,20} \\[4pt]
\text{VPD:} \quad VPD &= 0.75\big[6.1078\,e^{17.269 T_{\max}/(237.3+T_{\max})} - 6.1078\,e^{17.269 T_{\min}/(237.3+T_{\min})}\big] \tag{Eq\,88} \\[4pt]
\text{Water stress (photo):} \quad f_{w,\text{photo}} &= W_u / W_d \tag{Eq\,95}
\end{align}

\noindent In our proxy implementation, $f_{w,\text{photo}} \approx \min(\text{FASW}/0.5,\; 1)$ using gridded soil moisture divided by PAWC as a substitute for the full APSIM layered water balance.

% ======================================================================
\section{Notes}

\begin{itemize}
    \item \textbf{Yield units:} Wheat yield is converted from kg/ha to t/ha during extraction.
    \item \textbf{Missing values:} Cells with \texttt{no\_sow = 1} have zeros in all post-sowing features and NaN in yield. Filter these before ML training.
    \item \textbf{NaN in slot features:} \texttt{tmax\_max} and \texttt{tmin\_min} are NaN when a slot has zero valid days. Other variables are zero.
    \item \textbf{Temporal overlap:} Weekly and monthly features for the same period are computed independently --- they are not nested. XGBoost can select the most informative scale.
    \item \textbf{Pre-sowing features:} These capture antecedent conditions (stored soil water, autumn break rainfall) that influence sowing decisions and early crop establishment.
    \item \textbf{Configuration:} All parameters (paths, year range, thermal time windows, number of weekly/monthly slots) are in \texttt{config.py}. Temporal settings can be changed without modifying the extraction code.
\end{itemize}

\end{document}
